\par\Needspace{4\baselineskip}
\begin{center}
\SourceLeaf{031b}\pdfbookmark[2]{14.}{article-014}14.
\end{center}
Die Gleichungen in Art. 8 liefern mehrere Verfahren, um die wahre Anomalie und den Radiusvektor aus der exzentrischen Anomalie zu bestimmen. Die zweckmäßigsten davon werden wir erläutern.

\SourceLeaf{032}I. Beim üblichen Verfahren bestimmt man $v$ nach Gl. VII und anschließend $r$ nach Gl. II. Behält man für $p$ den in Art. 10 angegebenen Wert bei, so ergibt sich für das Beispiel des vorigen Artikels folgende Rechnung.
\begin{center}
\begin{tabular}{@{}ll@{}}
\multicolumn{2}{l}{$\frac{1}{2}E=162^0\,8^{\prime}\,14^{\prime\prime}{,}75$}\\[.3em]
\makebox[16em][l]{$\log\operatorname{tang}\frac{1}{2}E$\hfill .\,.\,.\,.\,.\,.} & $9{,}508\,2198_n$\\[.3em]
\makebox[16em][l]{$\log\operatorname{tang}(45^0-\frac{1}{2}\varphi)$\hfill .\,.} & $9{,}891\,2427$\\[.3em]
\hline\noalign{\vskip2pt}
\makebox[16em][l]{$\log\operatorname{tang}\frac{1}{2}v$\hfill .\,.\,.\,.\,.\,.\,.} & $9{,}616\,9771_n$\\[.3em]
\multicolumn{2}{l}{$\frac{1}{2}v=157^0\,30^{\prime}\,41^{\prime\prime}{,}50$}\\[.3em]
\multicolumn{2}{l}{$v=315\ \ 1\ \ 23{,}00$}\\[.3em]
\end{tabular}
\end{center}
\begin{center}
\begin{tabular}{@{}ll@{}}
\makebox[16em][l]{$\log e$\hfill .\,.\,.\,.\,.\,.\,.\,.\,.\,.\,.\,.} & $9{,}389\,7262$\\[.3em]
\makebox[16em][l]{$\log\cos v$\hfill .\,.\,.\,.\,.\,.\,.\,.\,.\,.} & $9{,}849\,6597$\\[.3em]
\hline
 & $9{,}239\,3859$\\[.3em]
$e\cos v=$ & $0{,}173\,5345$\\[.3em]
\hline
\makebox[16em][l]{$\log p$\hfill .\,.\,.\,.\,.\,.\,.\,.\,.\,.\,.\,.} & $0{,}395\,4837$\\[.3em]
\makebox[16em][l]{$\log(1+e\cos v)$\hfill .\,.\,.\,.\,.} & $0{,}069\,4959$\\[.3em]
\hline
\makebox[16em][l]{$\log r$\hfill .\,.\,.\,.\,.\,.\,.\,.\,.\,.\,.\,.} & $0{,}325\,9878$\\[.3em]
\end{tabular}
\end{center}

II. Das folgende Verfahren ist kürzer, wenn mehrere Positionen zu berechnen sind, da die konstanten Logarithmen der Größen $\BareRoot{(a(1+e))}$ und $\BareRoot{(a(1-e))}$ nur einmal berechnet werden müssen. Aus Gl. V und VI folgt
\[
\begin{gathered}
\sin\frac{1}{2}v\cdot\BareRoot{r}=\sin\frac{1}{2}E\BareRoot{(a(1+e))}\\[.5em]
\cos\frac{1}{2}v\cdot\BareRoot{r}=\cos\frac{1}{2}E\BareRoot{(a(1-e))},
\end{gathered}
\]
woraus sich $\frac{1}{2}v$ und $\log\BareRoot{r}$ leicht bestimmen lassen. Allgemein erhält man, wenn $P\sin Q=A$ und $P\cos Q=B$ gegeben sind, zunächst $Q$ aus der Formel $\operatorname{tang}Q=\frac{A}{B}$ und dann $P$ aus $P=\frac{A}{\sin Q}$ oder aus $P=\frac{B}{\cos Q}$. Die erste Formel ist vorzuziehen, wenn $\sin Q$ größer als $\cos Q$ ist, die zweite, wenn $\cos Q$ größer als $\sin Q$ ist. Aufgaben, die zu solchen Gleichungen führen (sie werden in diesem Werk sehr häufig vorkommen), enthalten meist die Bedingung, dass $P$ positiv sein muss. Damit ist unmittelbar entschieden, ob $Q$ zwischen $0$ und $180^0$ oder zwischen $180^0$ und $360^0$ zu wählen ist. Fehlt eine solche Bedingung, so bleibt uns diese Wahl überlassen.

In unserem Beispiel ist $e=0{,}245\,3162$
\begin{center}
\begin{tabular}{@{}ll@{}}
\makebox[16em][l]{$\log\sin\frac{1}{2}E$\hfill .\,.\,.\,.\,.\,.} & $9{,}486\,7632$\\[.3em]
\makebox[16em][l]{$\log\BareRoot{(a(1+e))}$\hfill .\,.\,.} & $0{,}258\,8593$\\[.3em]
\end{tabular}
\end{center}
\begin{center}
\begin{tabular}{@{}ll@{}}
\makebox[16em][l]{$\log\cos\frac{1}{2}E$\hfill .\,.\,.\,.\,.\,.} & $9{,}978\,5434_n$\\[.3em]
\makebox[16em][l]{$\log\BareRoot{(a(1-e))}$\hfill .\,.\,.} & $0{,}150\,1020$\\[.3em]
\end{tabular}
\end{center}

\SourceLeaf{033}\noindent Daraus folgt
\begin{center}
\footnotesize
\begin{tabular}{@{}ll@{}}
$\left.\begin{array}{ll}
\log\sin\frac{1}{2}v\cdot\BareRoot{r}\quad .\,. &9{,}745\,6225\\[.3em]
\log\cos\frac{1}{2}v\cdot\BareRoot{r}\quad .\,. &0{,}128\,6454_n
\end{array}\right\}$ & \begin{tabular}{@{}l@{}}
daraus $\log\operatorname{tang}\frac{1}{2}v$ .\,.\,.\,. $9{,}616\,9771_n$\\[.3em]
$\frac{1}{2}v=157^0\,30^{\prime}\,41^{\prime\prime}{,}50$\\[.3em]
$v=315\ \ 1\ \ 23{,}00$
\end{tabular}\\
\begin{tabular}{@{}ll@{}}
\makebox[13em][l]{$\log\cos\frac{1}{2}v$\hfill .\,.\,.\,.\,.} & $9{,}965\,6515_n$\\[.3em]
\hline
\makebox[13em][l]{$\log\BareRoot{r}$\hfill .\,.\,.\,.\,.\,.\,.} & $0{,}162\,9939$\\[.3em]
\makebox[13em][l]{$\log r$\hfill .\,.\,.\,.\,.\,.\,.\,.} & $0{,}325\,9878$\\[.3em]
\end{tabular} &
\end{tabular}
\end{center}

III. Wir fügen ein drittes Verfahren hinzu, das fast ebenso einfach durchzuführen ist wie das zweite. Für Rechnungen, die höchste Genauigkeit erfordern, ist es wegen seiner größeren Präzision meist vorzuziehen. Zunächst bestimmt man $r$ nach Gl. III und dann $v$ nach X. Damit ergibt sich für unser Beispiel folgende Rechnung.
\begin{center}
\begin{tabular}{@{}ll@{}}
\makebox[16em][l]{$\log e$\hfill .\,.\,.\,.\,.\,.\,.\,.\,.\,.\,.\,.} & $9{,}389\,7262$\\[.3em]
\makebox[16em][l]{$\log\cos E$\hfill .\,.\,.\,.\,.\,.\,.\,.\,.} & $9{,}909\,4637$\\[.3em]
\hline
 & $9{,}299\,1899$\\[.3em]
$e\cos E=$ & $0{,}199\,1544$\\[.3em]
\hline
\makebox[16em][l]{$\log a$\hfill .\,.\,.\,.\,.\,.\,.\,.\,.\,.\,.\,.} & $0{,}422\,4389$\\[.3em]
\makebox[16em][l]{$\log(1-e\cos E)$\hfill .\,.\,.\,.} & $9{,}903\,5488$\\[.3em]
\hline
\makebox[16em][l]{$\log r$\hfill .\,.\,.\,.\,.\,.\,.\,.\,.\,.\,.\,.} & $0{,}325\,9877$\\[.3em]
\end{tabular}
\end{center}
\begin{center}
\begin{tabular}{@{}ll@{}}
\makebox[16em][l]{$\log\sin E$\hfill .\,.\,.\,.\,.\,.\,.\,.\,.} & $9{,}766\,3366_n$\\[.3em]
\makebox[16em][l]{$\log\BareRoot{(1-e\cos E)}$\hfill .\,.\,.} & $9{,}951\,7744$\\[.3em]
\hline
 & $9{,}814\,5622_n$\\[.3em]
\makebox[16em][l]{$\log\sin\frac{1}{2}\varphi$\hfill .\,.\,.\,.\,.\,.\,.\,.} & $9{,}092\,0395$\\[.3em]
\hline\noalign{\vskip2pt}
\makebox[16em][l]{$\log\sin\frac{1}{2}(v-E)$\hfill .\,.\,.\,.} & $8{,}906\,6017_n$\\[.3em]
\multicolumn{2}{l}{$\frac{1}{2}(v-E)=-4^0\,37^{\prime}\,33^{\prime\prime}{,}24$}\\[.3em]
\multicolumn{2}{l}{$v-E=-9\ \ 15\ \ 6{,}48$}\\[.3em]
\multicolumn{2}{l}{$v=315\ \ 1\ \ 23{,}02$}\\[.3em]
\end{tabular}
\end{center}

Zur Kontrolle der Rechnung eignet sich Formel VIII oder XI sehr gut, besonders wenn $v$ und $r$ nach dem dritten Verfahren bestimmt wurden. Die Rechnung lautet wie folgt.
\begin{center}
\begin{tabular}{@{}ll@{}}
\makebox[16em][l]{$\log\frac{a}{r}\sin E$\hfill .\,.\,.\,.\,.} & $9{,}862\,7878_n$\\[.3em]
\makebox[16em][l]{$\log\cos\varphi$\hfill .\,.\,.\,.\,.\,.\,.} & $9{,}986\,5224$\\[.3em]
\hline
 & $9{,}849\,3102_n$\\[.3em]
\makebox[16em][l]{$\log\sin v$\hfill .\,.\,.\,.\,.\,.\,.} & $9{,}849\,31\overset{\scriptscriptstyle\text{‘}}{0}2_n$\\[.3em]
\end{tabular}
\end{center}
\begin{center}
\begin{tabular}{@{}ll@{}}
\makebox[16em][l]{$\log\sin E\BareRoot{\frac{a}{r}}$\hfill .\,.\,.\,.\,.} & $9{,}814\,5622_n$\\[.3em]
\makebox[16em][l]{$\log\cos\frac{1}{2}\varphi$\hfill .\,.\,.\,.\,.\,.} & $9{,}996\,6567$\\[.3em]
\hline
 & $9{,}811\,2189_n$\\[.3em]
\makebox[16em][l]{$\log\sin\frac{1}{2}(v+E)$\hfill .\,.\,.} & $9{,}811\,2189_n$\\[.3em]
\end{tabular}
\end{center}
